\chapterhead{4}{Muddy Water Through Paper}{Catching what is too big to pass}\label{ch4}

Muddy water is a mixture of water and tiny bits of soil. A sieve is no good, because the bits are too small. A \textbf{filter} is paper with very tiny holes. Water passes through the holes. The soil bits are too big and are caught.

\begin{center}
\begin{tikzpicture}
  \draw[line width=1.4pt,color=ink] (0,1.4) -- (1.0,0.2) -- (2.0,1.4);
  \draw[line width=1.4pt,color=ink] (-0.2,1.4) -- (2.2,1.4);
  \fill[brown!60!black!50] (0.55,0.5) -- (1.0,0.2) -- (1.45,0.5) -- cycle;
  \draw[-{Stealth},line width=1.6pt,color=sky!80!black] (1.0,0.1) -- (1.0,-0.6);
  \draw[line width=1.4pt,color=ink] (0.2,-0.7) -- (0.3,-1.8) -- (1.7,-1.8) -- (1.8,-0.7);
  \fill[sky!40] (0.33,-1.77) rectangle (1.67,-1.2);
  \node[font=\sffamily\scriptsize,anchor=west] at (2.6,0.9) {muddy water goes in};
  \node[font=\sffamily\scriptsize,anchor=west] at (2.6,0.3) {mud is caught on the paper};
  \node[font=\sffamily\scriptsize,anchor=west] at (2.6,-1.2) {clearer water comes out};
\end{tikzpicture}
\end{center}

\begin{tryit}{Filter muddy water}
\textbf{You need:} a kitchen scale, a coffee filter or paper towel, two cups, a spoon, a little garden soil, water.
\begin{enumerate}[leftmargin=1.6em]
\item Weigh a dry filter. Write it.
\item Stir about 10~g of soil into 190~g of water in a cup, to make 200~g of muddy water. Write 200.
\item Hold the filter over a clean cup. Pour the muddy water slowly into the filter.
\item Weigh the clear water collected. Then weigh the wet filter with the mud on it.
\end{enumerate}
\par\medskip
\renewcommand{\arraystretch}{1.6}
\begin{tabular}{|l|c|}
\hline
\textbf{Muddy water at the start (g)} & \\ \hline
\textbf{Dry filter (g)} & \\ \hline
\textbf{Clear water collected (g)} & \\ \hline
\textbf{Wet filter with mud (g)} & \\ \hline
\textbf{Mud and water on the filter (wet minus dry) (g)} & \\ \hline
\end{tabular}
\end{tryit}
\adult{Use garden soil that is free of sharp bits. Wash hands afterwards. Do not drink the water.}

\begin{worked}{Practice numbers, not from a real test}
Muddy water: 200~g. Dry filter: 2~g. Clear water collected: 186~g. Wet filter with mud: 14~g.\\
On the filter: $14-2=12$~g of mud and water.\\
Add up: $186+12=198$~g. The gap is $200-198=2$~g. Some water stayed in the cup and on the sides.
\end{worked}

\begin{center}
\begin{tikzpicture}
  \node[bigcir,minimum size=2.4cm] (a) at (0,0) {}; \node[font=\sffamily\scriptsize,align=center] at (0,0) {muddy\\water\\200};
  \node[cir] (b) at (4.6,1.1) {186}; \node[cir] (c) at (4.6,0) {12}; \node[cir] (d) at (4.6,-1.1) {2};
  \draw[popA] (a) -- (b); \draw[popA] (a) -- (c); \draw[popA] (a) -- (d);
  \node[note,anchor=west] at (5.3,1.1) {clearer water};
  \node[note,anchor=west] at (5.3,0) {mud on the filter};
  \node[note,anchor=west] at (5.3,-1.1) {the gap, named};
\end{tikzpicture}
\end{center}

\begin{multix}[hmm]\multixbot[hmm]\tcblower
The water coming through looks clean, so the filter must have made it pure. It looks clear to me. I like being sure. Should I be?
\end{multix}

\trust{We weighed the muddy water, the clear water and the wet filter, and named the small gap.}
{The three amounts added to within 2~g of the start, so the mud really stayed on the filter rather than disappearing.}
{That the clear water is pure. A filter catches what is too big for its holes. Dissolved salt is tiny and would pass straight through.}

\begin{keepbox}
We filtered muddy water. \textbf{We kept} \wline[3.4cm] \textbf{and we set aside} \wline[3.4cm].
\end{keepbox}

\clearpage

\subhead{Does the filter catch salt?}
\begin{tryit}{Muddy salt water}
Dissolve 5~g of salt in 100~g of water, add 5~g of soil, stir and filter. Put 25~g of the filtered water on a weighed saucer and let it dry. Is there a crust? Weigh it.\\[4pt]
Empty saucer \wline[2cm]~g \quad Dry saucer \wline[2cm]~g \quad Crust \wline[2cm]~g
\end{tryit}

\begin{worked}{Practice numbers}
Empty saucer: 30~g. Dry saucer: 33~g. Crust: 3~g. The filtered water was clear, yet a crust of salt remained. The salt had passed through the filter.
\end{worked}

\subhead{Which tool for which difference?}
\begin{center}
\sffamily\small\renewcommand{\arraystretch}{1.5}
\begin{tabular}{@{}p{3.6cm}p{3.5cm}p{3.5cm}@{}}
\toprule
\textbf{Tool} & \textbf{Catches} & \textbf{Lets through}\\
\midrule
Sieve & big bits & small bits and water\\
Filter paper & soil, sand, bits you can see or cloud & water with dissolved salt\\
Evaporation & nothing (water leaves) & salt stays behind\\
\bottomrule
\end{tabular}
\end{center}

\clearpage

\begin{puzzle}{Add up the filter}
Muddy water 200~g. Clear water 185~g. Mud and water on the filter 13~g. Total \wline[1.6cm]~g. Gap \wline[1.6cm]~g.
\end{puzzle}

\begin{puzzle}{Wet minus dry}
A dry filter weighs 2~g. After filtering it weighs 17~g. How much mud and water is on it? \wline[1.6cm]~g
\end{puzzle}

\begin{puzzle}{Choose the tool}
Write S for sieve, F for filter, E for evaporation or X if none will work.\\
(a) sand from water \wline[1cm] \quad (b) salt from clear salt water \wline[1cm] \quad (c) pebbles from sand \wline[1cm]
\end{puzzle}

\begin{puzzle}{Clear but not pure}
A filter turned muddy water clear. A student says it is now pure. How could she test? \wline[6cm]
\end{puzzle}

\begin{puzzle}{Spot the unfair test}
Tom filtered muddy water through a coffee filter. Lena filtered it through a paper towel. Lena's water was cloudier, so she says muddy water cannot be filtered. What is wrong? \wline[6cm]
\end{puzzle}

\clearpage

\begin{puzzle}{Draw the circles}
Muddy water of 150~g gave 140~g of clearer water and 8~g on the filter. Fill in the gap.
\begin{center}
\begin{tikzpicture}
  \node[cir] (a) at (0,0) {150}; \node[cir] (b) at (4,0.9) {140}; \node[cir] (c) at (4,0) {8}; \node[cir,dashed] (d) at (4,-0.9) {};
  \draw[popA] (a) -- (b); \draw[popA] (a) -- (c); \draw[popA] (a) -- (d);
\end{tikzpicture}
\end{center}
\end{puzzle}

\begin{puzzle}{Is the gap too big?}
The gap in the last puzzle is \wline[1.6cm]~g. Is that within the scale's 1~g error? \wline[2cm] What would you check? \wline[4cm]
\end{puzzle}

\begin{puzzle}{Sentence check}
Mark S, G or B.\\
(a) ``The filtered water weighed 186~g.'' \wline[1cm]\\
(b) ``The filter caught all the soil.'' \wline[1cm]\\
(c) ``Filtering makes any water safe.'' \wline[1cm]
\end{puzzle}

\begin{puzzle}{Cloudy water}
A cup of cloudy water was left to stand overnight. The next day the bottom of the cup was muddy. Explain. \wline[6cm]
\end{puzzle}

\begin{puzzle}{Your test}
Plan a fair test of which paper filters best: coffee filter, paper towel or tissue. Change: \wline[3cm] Keep the same: \wline[4cm] Measure: \wline[3cm]
\end{puzzle}
\clearpage

\thispagestyle{fancy}
\subhead{Filter comparison}
Filter the same amount of the same muddy water through two or three different papers. Keep everything else the same.\par\medskip
\renewcommand{\arraystretch}{3.0}
\begin{tabularx}{\linewidth}{|>{\sffamily\small\bfseries}p{2.8cm}|X|X|X|X|}
\hline
Paper & Muddy water in (g) & Clear water out (g) & Left on paper (g) & How clear? \\ \hline
 & & & & \\ \hline
 & & & & \\ \hline
 & & & & \\ \hline
\end{tabularx}

\vspace{8pt}
Which paper was clearest? \wline[6cm]\\[10pt]
Which was quickest? \wline[6cm]\\[10pt]
Is the quickest also the clearest? Say what that tells you. \wline[4cm]\\[10pt]
Draw the cups side by side to show how clear each was.\par\medskip
\begin{tikzpicture}\draw[ink!70,line width=1pt,rounded corners=3pt] (0,0) rectangle (11.4,4.2);\end{tikzpicture}
\clearpage
